\section{Related Work}

\subsection{Anisotropy in Representation Geometry}

Anisotropy in learned representations has been widely observed in both static
and contextual embedding models.
Ethayarajh~\cite{ethayarajh2019contextual} showed that contextualized
representations from BERT, ELMo, and GPT-2 occupy a narrow cone in
representation space, exhibiting strong directional concentration.
Mu and Viswanath~\cite{mu2018all,mu2017allbutthetop} demonstrated that
removing the top principal components from word embeddings significantly
improves isotropy and downstream semantic similarity.
Arora et al.~\cite{arora2017simple} similarly proposed subtracting a
common discourse vector in sentence embeddings to mitigate dominant
shared components.
Subsequent analyses identified “rogue dimensions” and dominant outlier
directions in Transformer models
\cite{timkey2021rogue,kovaleva2021bertbusters},
while more recent work argued that anisotropy may depend on training
conditions rather than being inherent to Transformer architecture
\cite{machina2024anisotropy}.

These studies establish anisotropy as a persistent geometric phenomenon.
However, they primarily treat it as a spectral artifact to be post-processed
or regularized.
In contrast, our work identifies a specific structural source of anisotropy:
a coherent activation mean bias originating from the embedding layer,
which propagates through residual connections and dominates elementwise
extremes during training.

\subsection{Frequency Effects and Embedding Structure}

The relationship between word frequency and embedding geometry has been
recognized since early distributed representation models.
Word2Vec~\cite{mikolov2013efficient,mikolov2013distributed}
introduced subsampling of frequent tokens to reduce dominance of high-frequency
words during training.
Arora et al.~\cite{arora2017simple} provided a generative model explaining
why common words induce a shared background direction in embeddings.
Mu and Viswanath~\cite{mu2018all} formalized the removal of dominant principal
components to restore isotropy.
More recently, Puccetti et al.~\cite{puccetti2022frequency}
argued that outlier dimensions in Transformers are driven by token frequency.

While prior work links frequency to dominant directions,
it largely focuses on representational quality.
We extend this perspective to training dynamics,
showing that frequency-induced embedding bias becomes the dominant driver
of activation outliers and optimizer second-moment inflation,
with direct implications for large-scale training stability.

\subsection{Outlier Suppression and Low-Bit Quantization}

Outlier dimensions and extreme activation magnitudes have been a central
challenge in low-bit quantization of large language models.
SmoothQuant~\cite{xiao2023smoothquant} redistributes activation and weight
scales to reduce outlier concentration.
GPTQ~\cite{frantar2022gptq}, AWQ~\cite{lin2024awq}, QLoRA~\cite{dettmers2023qlora},
and SVDQuant~\cite{li2024svdquant}
propose various strategies to absorb or compensate for large-magnitude
components.
QuaRot~\cite{ashkboos2024quarot} and HALO~\cite{ashkboos2025halo}
employ orthogonal rotations or Hadamard transforms to mitigate outliers.
Outlier Suppression~\cite{wei2023outlier} explicitly targets extreme
activation dimensions.
Recent work on FP8 and FP4 training further highlights the sensitivity of
low-precision regimes to activation dynamic range
\cite{micikevicius2022fp8,peng2023fp8,perez2023fp8,cao2025metis}.

These methods typically treat outliers as spectral imbalances requiring
rotation, clipping, or decomposition.
Our analysis shows that a large fraction of extreme activation magnitudes
is attributable to a rank-one mean bias component.
Addressing this bias at the embedding activation level
significantly reduces dynamic-range inflation before more complex spectral
corrections are applied.

\subsection{Efficient Spectral Approximation and Whitening}

Low-rank approximation and spectral methods are classical tools for
controlling anisotropy.
Randomized SVD~\cite{halko2011randomized} and distributed low-rank
algorithms~\cite{suresh2017distributed}
provide scalable approximations of dominant singular directions.
Oja’s rule~\cite{oja1982pca} and whitening-based approaches
\cite{huang2021whiteningbert}
offer incremental spectral normalization.

However, applying SVD or whitening repeatedly across large models
remains computationally expensive.
Our work shows that a significant portion of harmful anisotropy
is concentrated in a coherent rank-one bias component,
allowing a single mean-subtraction operation to achieve much of the
desired conditioning effect at minimal cost.
Residual low-rank whitening can then be applied selectively at the
embedding activation level rather than throughout the network.

\paragraph{Positioning.}
Prior work has identified anisotropy, frequency-driven dominant directions,
and outlier suppression strategies.
We synthesize these threads by demonstrating that
mean-dominated activation bias constitutes a primary structural source
of anisotropy in large language model training.
By tracing this bias to embedding frequency statistics and quantifying its
impact on both quantization and optimization,
we provide a unified perspective that connects representation geometry,
training stability, and hardware efficiency.